\documentclass[a4paper]{article}

% Set margins
\usepackage[hmargin=2.5cm, vmargin=3cm]{geometry}

\frenchspacing

% Language packages
\usepackage[utf8]{inputenc}
\usepackage[T1]{fontenc}
\usepackage[magyar]{babel}

% AMS
\usepackage{amssymb,amsmath}

% Graphic packages
\usepackage{graphicx}

% Colors
\usepackage{color}
\usepackage[usenames,dvipsnames]{xcolor}

% Listings
\usepackage{listings}

% Question
\newenvironment{question}[1]
{\noindent\textcolor{OliveGreen}{$\circ$ \textit{#1}}

\smallskip

\color{Gray}

}{\bigskip}

% Task
\newenvironment{task}[1]
{\noindent\textcolor{RoyalBlue}{$\circ$ \textit{#1}}

\smallskip

\color{Gray}

}{\bigskip}

% Notification
\newenvironment{notification}[1]
{\noindent\textcolor{Peach}{$\circ$ \textit{#1}}

\smallskip

\color{Gray}

}{\bigskip}

% Problem
\newenvironment{problem}[1]
{\noindent\textcolor{OrangeRed}{$\circ$ \textit{#1}}

\smallskip

\color{Gray}

}{\bigskip}

% Solution
\newenvironment{solution}
{\color{RoyalBlue}}{\bigskip}

% Comment
\newenvironment{comment}
{\color{Green}}{\bigskip}

\lstset{% setup listings
        framexleftmargin=16pt,
        framextopmargin=6pt,
        framexbottommargin=6pt, 
        frame=single, rulecolor=\color{black}
}

\begin{document}

\begin{center}
   \large \textbf{Valószínűségszámítás - típusfeladatok}
\end{center}

\vskip 1cm

\section{Típusfeladatok}

\begin{itemize}
\item A grafikus ábrázolás előnyös, de érdemes minél részletesebben, áttekinthető formában számításokkal ki is egészíteni.
\item A szöveges feladatokban a jelölések nem kötöttek, ezért azokat mindig külön definiálni kell. (Például, hogy egy $A$ esemény mit jelent.)
\item Az eloszlás- vagy sűrűségfüggvény felrajzolására is szükség lehet.
\item Bizonyos feladatoknál vannak hiányzó adatok, amelyeket az ismert összefüggések alapján kell meghatározni.
\item A függvények ábrázolásánál figyelni kell arra, hogy minden tengelyfelirat, lépték, intervallum (hol nyitott, hol zárt) egyértelműen látszódjon.
Ha a grafikon nehezen áttekinthető, vagy javítani kellett benne, akkor érdemes lehet újrarajzolni.
\item Figyelni kell arra, hogy a feladatban szórás vagy szórásnégyzet meghatározása szerepel-e!
\item Mivel többségében szöveges feladatokról van szó, ezért a kapott eredményeket szöveges formában is össze kell foglalni. (Tehát ténylegesen úgy, ahogy a kérdésre válaszolnánk.)
\end{itemize}

\subsection{Klasszikus valószínűségi mező}

A szöveges leírás alapján a kedvező és az összes eset arányát kell meghatározni.

\subsection{Geometriai valószínűség}

A feladat szöveges formában adott, amelyet valamilyen geometriai valószínűségre kell visszavezetni.

\subsection{Visszatevéses, visszatevés nélküli mintavétel}

A képletekkel számolható.

\subsection{Feltételes valószínűség, átrendezés}

Az eseményekhez meg van adva néhány valószínűség és feltételes valószínűség.
Ezek alapján kell a hiányzó valószínűségeket meghatározni.

\subsection{Feltételes valószínűség, Bayes tétel}

A feladat felírható táblázat formájában, amelyben vagy a gyakoriságok vagy a valószínűségek szerepelnek.
Az adott eseményekhez tartozó valószínűségeket, feltételes valószínűségeket kell meghatározni.

\subsection{Diszkrét eloszlás}

Meg van adva az eloszlás felvett értékekkel és valószínűségekkel.
Meg kell határozni a várható értéket és a szórást.
Fel kell tudni rajzolni az eloszlásfüggvényt.

\subsection{Folytonos eloszlás}

Az eloszlás adott a sűrűség- vagy az eloszlásfüggvényével.
Tipikusan hiányozhatnak belőle paraméterek (például konstans vagy intervallum vége).
\begin{itemize}
\item Meg kell határozni a hiányzó értékeket.
\item Ki kell számítani a várható értéket és a szórást.
\item Meg kell határozni adott intervallumba esésnek a valószínűségét.
\item A függvényeket szükséges lehet ábrázolni.
\end{itemize}

\subsection{Eseménysor}

Egy kísérlet ismétlődik valamilyen feltétel teljesüléséig.
Meg van adva közvetve, vagy közvetlenül, hogy mennyi az adott esemény bekövetkezésének a valószínűsége.
Meg kell határozni az eseménysor várható hosszát, a hossz eloszlásának szórását, a bekövetkezés valószínűségét bizonyos intervallumon.

\subsection{Egyenletes eloszlás}

$U[a, b]$, várható értéket, szórást, intervallumba esést kell számolni, a sűrűségfüggvényt és az eloszlásfüggvényt kell tudni ábrázolni.

\subsection{Binomiális eloszlás}

Adott a $p$ valószínűség, és az $n$, amely a végrehajtott próbálkozások/mintavételek számát írja le.
Meg kell határozni a várható értéket, a szórást, és adott intervallumon belül, hogy mennyi a valószínűség.

\end{document}
